\documentclass[tikz,border=8pt]{standalone}
\usepackage[no-math]{fontspec}
\usepackage[brazil]{babel}
\usepackage{amsmath,amssymb}
\IfFileExists{C:/Windows/Fonts/Lato_Regular.ttf}{%
  \setmainfont{Lato}[
    Path = C:/Windows/Fonts/,
    UprightFont = Lato_Regular.ttf,
    BoldFont = Lato_Bold.ttf,
    ItalicFont = Lato_Italic.ttf,
    BoldItalicFont = Lato_BoldItalic.ttf
  ]
  \setsansfont{Lato}[
    Path = C:/Windows/Fonts/,
    UprightFont = Lato_Regular.ttf,
    BoldFont = Lato_Bold.ttf,
    ItalicFont = Lato_Italic.ttf,
    BoldItalicFont = Lato_BoldItalic.ttf
  ]
}{%
  \RequirePackage[default]{lato}
}
\renewcommand{\familydefault}{\sfdefault}

\usepackage{xcolor}
\usepackage{tikz}
\usepackage{pgfplots}
\pgfplotsset{compat=1.18}
\usetikzlibrary{arrows.meta,positioning,shapes.geometric,fit,backgrounds,calc,decorations.pathreplacing,patterns}

% Paleta institucional UEFS / CHF662
\definecolor{uefsNavy}{HTML}{1B2A4A}
\definecolor{uefsTeal}{HTML}{0E7C7B}
\definecolor{uefsForest}{HTML}{2D6A4F}
\definecolor{uefsGold}{HTML}{C4983B}
\definecolor{uefsGoldDark}{HTML}{8F6B1E}
\definecolor{uefsAlert}{HTML}{B02A37}
\definecolor{uefsOrange}{HTML}{D95F02}
\definecolor{uefsInk}{HTML}{22252A}
\definecolor{uefsMuted}{HTML}{5C6470}
\definecolor{uefsWater}{HTML}{2E86AB}
\definecolor{uefsSand}{HTML}{D9B382}
\definecolor{bgNavy}{HTML}{F2F5F9}
\definecolor{bgTeal}{HTML}{E9F5F4}
\definecolor{bgForest}{HTML}{EEF5F0}
\definecolor{bgGold}{HTML}{FCF7EB}
\definecolor{bgAlert}{HTML}{FDF2F2}
\definecolor{bgGray}{HTML}{F8FAFC}
\definecolor{lineGray}{HTML}{CFD8DC}

\tikzset{
  >=Stealth,
  aaCard/.style={draw=#1, fill=white, rounded corners=4pt, line width=1pt, inner sep=6pt, align=center, font=\small},
  aaBox/.style={draw=#1, rounded corners=4pt, line width=1.1pt, inner sep=6pt, align=center, font=\footnotesize\sffamily, text=uefsInk},
  aaTitle/.style={font=\normalsize\bfseries\sffamily, text=uefsNavy, align=center},
  aaSub/.style={font=\footnotesize\sffamily, text=uefsMuted, align=center},
  aaArrow/.style={->, >=Stealth, line width=1.4pt},
  aaDArrow/.style={<->, >=Stealth, line width=1.4pt}
}


\begin{document}
\begin{tikzpicture}[font=\sffamily]

% Título principal
\node[font=\large\bfseries\sffamily, text=uefsNavy, align=center] at (0, 3.85)
  {Suscetibilidade, Fragilidade, Vulnerabilidade e Risco Socioambiental};

% =========================================================
% LINHA SUPERIOR: 3 DIMENSÕES
% =========================================================
\node[draw=uefsAlert, fill=bgAlert, rounded corners=6pt, line width=1.2pt,
      text width=4.95cm, minimum height=2.55cm, inner sep=8pt, align=center,
      font=\footnotesize\sffamily, text=uefsInk] (perigo) at (-5.75, 1.75)
  {\textbf{\small\color{uefsAlert}Perigo / Ameaça ($P$)\\[-1pt]e Suscetibilidade}\\[4pt]
   \textit{\color{uefsMuted}(dimensão físico-natural)}\\[4pt]
   Predisposição do meio e probabilidade de evento físico adverso.};

\node[draw=uefsGoldDark, fill=bgGold, rounded corners=6pt, line width=1.2pt,
      text width=4.95cm, minimum height=2.55cm, inner sep=8pt, align=center,
      font=\footnotesize\sffamily, text=uefsInk] (frag) at (0, 1.75)
  {\textbf{\small\color{uefsGoldDark}Fragilidade Ambiental\\[-1pt](Ross; Crepani)}\\[4pt]
   \textit{\color{uefsMuted}(potencial vs. emergente)}\\[4pt]
   Sensibilidade intrínseca do meio modulada pelo uso da terra.};

\node[draw=uefsWater, fill=bgNavy, rounded corners=6pt, line width=1.2pt,
      text width=4.95cm, minimum height=2.55cm, inner sep=8pt, align=center,
      font=\footnotesize\sffamily, text=uefsInk] (vuln) at (5.75, 1.75)
  {\textbf{\small\color{uefsNavy}Exposição ($E$) e\\[-1pt]Vulnerabilidade ($V$)}\\[4pt]
   \textit{\color{uefsMuted}(dimensão socioespacial)}\\[4pt]
   População e bens expostos somados à fragilidade social de resposta.};

% =========================================================
% SETAS VERTICAIS PARA O BANNER DE RISCO
% =========================================================
\draw[->, >=Stealth, line width=1.8pt, uefsAlert] (perigo.south) -- (-5.75, -0.25);
\draw[->, >=Stealth, line width=1.8pt, uefsGoldDark] (frag.south) -- (0, -0.25);
\draw[->, >=Stealth, line width=1.8pt, uefsNavy] (vuln.south) -- (5.75, -0.25);

% =========================================================
% BANNER CENTRAL ESCURO: FÓRMULA DO RISCO
% =========================================================
\node[draw=uefsGold, fill=uefsNavy, rounded corners=7pt, line width=1.5pt,
      text width=16.2cm, minimum height=1.55cm, inner sep=9pt, align=center,
      text=white] (risco) at (0, -1.1)
  {\normalsize\bfseries\color{uefsGold}Risco Socioambiental: \quad $R = P \times E \times V$\\[5pt]
   \footnotesize\normalfont\itshape\color{white}Sem exposição ($E=0$), há perigo natural, mas não há risco social $\cdot$ Graus IPT R1 a R4};

% =========================================================
% PÍLULA INFERIOR: CAPACIDADE DE SUPORTE
% =========================================================
\node[draw=uefsTeal, fill=bgTeal, rounded corners=11pt, line width=1.2pt,
      text width=16.2cm, minimum height=0.85cm, inner sep=7pt, align=center,
      font=\small\sffamily, text=uefsForest] at (0, -2.55)
  {\textbf{\color{uefsTeal}Capacidade de Suporte:} limite máximo de pressão suportado antes do colapso da resiliência.};

\end{tikzpicture}
\end{document}
